\documentclass[acmtog]{acmart}
\usepackage{graphicx}
\usepackage{subfigure}
\usepackage{natbib}
\usepackage{listings}
\usepackage{bm}
\usepackage{amsmath}

\definecolor{blve}{rgb}{0.3372549, 0.61176471, 0.83921569}
\definecolor{gr33n}{rgb}{0.29019608, 0.7372549, 0.64705882}
\makeatletter
\lst@InstallKeywords k{class}{classstyle}\slshape{classstyle}{}ld
\makeatother
\lstset{language=C++,
  basicstyle=\ttfamily,
  keywordstyle=\color{blve}\ttfamily,
  stringstyle=\color{red}\ttfamily,
  commentstyle=\color{magenta}\ttfamily,
  morecomment=[l][\color{magenta}]{\#},
  classstyle = \bfseries\color{gr33n},
  tabsize=2
}
\lstset{basicstyle=\ttfamily}

% Title portion
\title{Assignment 1:\\ {Exploring OpenGL Programming}}

\author{Name: Xu Minhao\\ student number: 2024533139}

% Document starts
\begin{document}
  \maketitle

  \vspace*{2 ex}

  \section{Introduction}
  This assignment implements a small interactive OpenGL application for rendering a
  three-dimensional mesh.  The program creates an OpenGL 3.3 core-profile
  context, loads a sphere from an object file, and renders it with depth testing
  and per-fragment lighting.  In addition to the basic rendering pipeline, the
  application supports first-person camera movement, mouse-look control, and
  keyboard rotation.  The main purpose of the assignment is to become familiar
  with the relationship between CPU-side OpenGL resource management, shader
  programs, vertex attributes, and the transformation pipeline.

  \section{Implementation Details}
  \subsection{Window and OpenGL Initialization}
  The \texttt{Window} class wraps GLFW window creation.  It requests an OpenGL
  3.3 core context, makes the context current, and initializes the OpenGL
  function pointers through GLAD.  The initial viewport is obtained from the
  framebuffer size, and a framebuffer-size callback updates the viewport when
  the window is resized.  Vertical synchronization is enabled through
  \texttt{glfwSwapInterval(1)}.

  \subsection{Mesh Loading and Vertex Specification}
  The \texttt{Mesh} class reads the vertex positions, vertex normals, and
  triangle face records from the supplied \texttt{.object} file.  Each face is
  expanded into three vertices, where every vertex stores a position and a
  normal.  The data are uploaded to a vertex buffer object (VBO), while an
  element buffer object (EBO) stores the indices of the triangles.  A vertex
  array object (VAO) records the vertex layout:
  attribute 0 contains three floating-point position values and attribute 1
  contains three floating-point normal values.  During drawing, the VAO is
  bound and \texttt{glDrawElements} renders the triangles.

  \subsection{Shader Pipeline}
  The vertex shader transforms each input position from model space to clip
  space:
  \[
    \boldsymbol{p}_{clip}
    = P V M
      \begin{bmatrix}\boldsymbol{p}_{model}\\1\end{bmatrix}.
  \]
  It also transforms the normal with the normal matrix and passes the world
  space position and normal to the fragment shader.  The fragment shader
  computes ambient, diffuse, and specular terms.  In particular, the diffuse
  term is based on the non-negative dot product between the normalized normal
  and light direction, while the specular term uses the reflected light
  direction and a shininess exponent of 32:
  \[
    I = k_a C +
        k_d \max(\boldsymbol{n}\cdot\boldsymbol{l},0)C +
        k_s \left(\max(\boldsymbol{v}\cdot\boldsymbol{r},0)\right)^{32}.
  \]
  The object color is orange-brown, and the light is placed at
  $(2,3,2)$ in world space.

  \subsection{Camera and Interaction}
  The camera stores its position, viewing direction, yaw, and pitch.  The view
  matrix is constructed from the camera's forward, right, and up vectors, which
  are obtained using vector cross products and normalization.  A perspective
  projection with a 45-degree vertical field of view, near plane 0.1, and far
  plane 100.0 is recomputed using the current framebuffer aspect ratio.

  The \texttt{W}, \texttt{A}, \texttt{S}, and \texttt{D} keys move the camera
  relative to its viewing direction, while \texttt{E} and \texttt{Q} move it
  along the world-up axis.  Movement is multiplied by frame time, so it is
  independent of the frame rate.  The left and right arrow keys change yaw,
  and the up and down arrow keys change pitch.  Mouse motion provides a second
  way to control yaw and pitch; pitch is clamped to the interval
  $[-89^\circ,89^\circ]$ to avoid the singularity when the camera looks
  directly along the up axis.  Pressing Escape closes the window.

  \subsection{Multiple Objects and Motion}
  The scene is not limited to a single mesh.  The program creates several
  \texttt{ScreenObject} entries, each storing a mesh pointer, a model matrix,
  a position, a velocity, a scale, a normal matrix, and an object color.  These
  entries are placed in a \texttt{std::array<ScreenObject, 3>} and rendered in a
  single loop.  This allows the same shader pipeline to be reused for multiple
  meshes, while still keeping each object's transform data independent.

  At every frame, the model matrix is recomputed from the object's current
  position and scale, and the normal matrix is updated accordingly.  The camera
  and light parameters are then uploaded to the shader as uniforms before the
  meshes are drawn.  A simple animation is also included by updating each
  object's position with velocity multiplied by the elapsed time:
  \[
    \mathbf{p}_{t+1} = \mathbf{p}_t + \mathbf{v}\,\Delta t.
  \]
  This makes the objects move smoothly across the scene while maintaining a
  consistent transformation pipeline for each rendered mesh.

  \section{Results}
  When executed, the application opens a $1280\times720$ GLFW window and
  renders the sphere against a blue-gray background.  The depth buffer is
  enabled, so nearer parts of the sphere correctly occlude farther parts.
  The lighting model produces a constant ambient contribution, a diffuse
  shading gradient across the surface, and a visible specular highlight near
  the reflection direction.  The rendered object remains correctly framed
  when the window is resized because both the viewport and the projection
  aspect ratio are updated.

  The interactive camera can orbit and translate around the object using the
  mouse and keyboard controls described above.  This also makes the effects of
  the view matrix and perspective projection directly observable.  Overall,
  the program demonstrates the complete path from mesh-file data, through
  OpenGL buffers and shader stages, to a depth-tested, interactively viewed
  image.

  \section{Build Diagnostics and Fix Note}
  During the compilation stage, the project initially failed in the object-update
  loop inside the main rendering routine.  The compiler reported that the
  assignment and compound-assignment operators were unavailable for the loop
  variable, indicating that the object was treated as a constant copy rather
  than a mutable element of the array.  In particular, the loop used
  \texttt{for (auto const object : objects)}, which prevents updates to
  \texttt{object.model}, \texttt{object.normalMatrix}, and
  \texttt{object.position}.  The correct pattern is to iterate by reference,
  \texttt{for (auto \&object : objects)}, so that the position, velocity, and
  model matrix can be recomputed each frame.  This fix resolves the reported
  \texttt{operator=} and \texttt{operator+=} errors and allows the animation and
  per-frame matrix updates to proceed normally.

  \section{Conclusion}
  This assignment provided practical experience with the core components of an
  OpenGL rendering pipeline.  The most important implementation details were
  preserving the VAO/VBO/EBO state relationships, passing data between shader
  stages with matching interfaces, and keeping the camera matrices and
  framebuffer aspect ratio consistent.  The final program combines these
  components into a compact interactive renderer with basic Phong-style
  illumination.

\end{document}
